% =====================================================================
% Ensaio 1 — Seção 7: Conclusão
%
% ⚠️ ESTA SEÇÃO TEM DE ENCARAR UM PROBLEMA QUE O PRÓPRIO TEXTO CRIOU.
%
% A §3 declara, em citação destacada: "Um Ensaio 1 que entregasse apenas o
% efeito médio no nível tratado deixaria o Ensaio 2 sem âncora empírica."
% A §6 entrega exatamente isso, e declara a derivada não-interpretável.
%
% Escrever uma conclusão que não diga isso seria escrever contra o próprio
% documento. A §7.2 diz, e propõe a saída.
%
% ⚠️ DECISÕES QUE SÃO DO PESQUISADOR, NÃO DESTE TEXTO: estão marcadas no corpo
% como tais, cada uma com recomendação. Não são perguntas retóricas.
% =====================================================================

\section{Conclusão}
\label{sec:conclusao}

\subsection{O que este ensaio estabelece}
\label{sub:estabelece}

Este ensaio montou um desenho de avaliação causal para o banimento da
pulverização aérea no Ceará e o levou até o fim: identificação fechada, dose
construída e verificada contra o dado, plano de análise datado e versionado,
estimação executada e camada de robustez inteira aplicada.

O que ele \textbf{não} entrega é efeito medido. O \(ATT(d\,|\,d)\) agregado é de
\(-36{,}19\) g com meia-largura de 30,32 g, contra o piso de 15 g fixado antes
de olhar, e nenhum dos três procedimentos de inferência rejeita o nulo. Pelo
critério que o próprio trabalho escreveu antes de ver qualquer coeficiente,
isso é \textbf{limite superior informativo}, e é assim que está reportado.

\begin{quote}
\textbf{E o resultado mais sólido deste ensaio é sobre o desenho, não sobre o
Ceará.} A análise de poder \emph{ex ante} previu efeito mínimo detectável de
33,4 g; a meia-largura realizada foi de 30,3 g. O desenho comportou-se como
diagnosticado antes de existir dado --- o que valida o procedimento e condena a
especificação corrente a não distinguir o efeito procurado do zero.
\end{quote}

\subsection{A âncora que o Ensaio 2 esperava, e que não veio}
\label{sub:ancora}

A §\ref{sec:teoria} justifica a escolha do parâmetro-alvo com um argumento que
agora se volta contra o trabalho, e é preciso dizê-lo sem rodeio. Aquela seção
afirma que Weitzman decide pela \emph{inclinação} e não pelo nível, e conclui
que um Ensaio 1 limitado ao efeito médio deixaria o Ensaio 2 sem âncora
empírica, devolvendo a escolha entre proibir e taxar ao terreno da
plausibilidade.

\textbf{É o que aconteceu.} A curva foi estimada, e o agregado que ela produz não
é interpretável: a derivada é fortemente não linear, inverte de sinal ao longo da
dose, e o agregado é média simples dominada pela faixa de dose quase nula, onde
o erro-padrão supera a estimativa. Nenhum dos 169 pontos exclui o zero, em
nenhuma faixa. A troca do alvo primário está registrada na tabela de desvios da
pré-especificação, mas o registro resolve a integridade do procedimento, não o
problema substantivo.

\paragraph{Três rotas, e a que se recomenda.}
A situação admite três saídas, e elas não são equivalentes.

\begin{enumerate}
\item \textbf{Adiar o Ensaio 2 até o Ensaio 1 ter poder.} A via declarada é
  encorpar o grupo tratado, baixando o corte de dose para além do decil
  superior. É honesta e é cara: serializa a dissertação inteira atrás de uma
  reestimação cujo ganho de precisão não está garantido.

\item \textbf{Substituir o ponto por limites.} \citet{callaway2024continuous}
  oferecem, na §5.1, identificação parcial sob hipótese de direção do viés de
  seleção. Entrega faixa, não número --- e uma faixa larga demais devolve o
  mesmo problema com mais aparato.

\item \textbf{Inverter a pergunta.} Em vez de estimar a inclinação e alimentar a
  regra de Weitzman, perguntar \textbf{qual inclinação seria necessária para
  inverter o ranking entre proibir e taxar}, e então confrontar esse valor com o
  que o Ensaio 1 consegue excluir.
\end{enumerate}

\begin{quote}
\textbf{Recomenda-se a terceira.} Ela converte uma estimativa que falhou em um
exercício de limites que é informativo mesmo com poder curto: o intervalo por
aleatorização, de \([-46{,}45;\ +57{,}85]\), \emph{exclui} regiões do espaço de
parâmetros, e se o valor que inverteria o ranking estiver fora do que o dado
admite, a conclusão de política se sustenta sem depender do ponto. É a forma
padrão de fazer análise de política com desenho imperfeito, e não exige nenhum
dado novo.
\end{quote}

\paragraph{Decisão do pesquisador.}
Qual das três rotas o Ensaio 2 adota é escolha de escopo da dissertação, e não
decorre do que este ensaio mediu. As três são defensáveis; a segunda e a
terceira são combináveis. O que \emph{não} é defensável é manter a redação atual
da §\ref{sec:teoria}, que promete uma âncora que o Ensaio 1 não entrega --- essa
promessa precisa ser reescrita na versão final, qualquer que seja a rota.

\subsection{Três achados que não são sobre o Ceará}
\label{sub:metodologicos}

O trabalho produziu resultados que sobrevivem independentemente do efeito do
banimento, e que valem para quem montar desenho parecido.

\paragraph{O agregado da curva pode não significar nada.}
Quando a densidade da dose se concentra perto de zero, a média simples das
derivadas pontuais é dominada pela região em que a derivada não é identificada.
O número sai, tem erro-padrão, e não quer dizer o que parece. Reportar a
decomposição por faixa de dose deveria ser rotina, e não é.

\paragraph{Fonte com cobertura excelente pode ser a mais perigosa.}
O SISAGUA cobre os 184 municípios em todos os anos da janela --- a melhor
cobertura de qualquer fonte deste trabalho --- e foi suspenso. As detecções de
agrotóxico no Ceará caem de 44 em 2019 para zero exato em 2020, 2021 e 2022,
\emph{com mais amostras coletadas}, enquanto o país segue reportando cerca de
metade dos resultados em valor numérico. Um diferença-em-diferenças sobre essa
variável concluiria que a lei eliminou a totalidade das detecções: grande,
significativo e inteiramente artefato de mudança de prática laboratorial. A
cobertura alta é o que torna o artefato convincente.

\paragraph{Um teste que ``passa'' precisa dizer com quanto poder.}
Medido formalmente, o placebo deste desenho exclui tendências confundidoras
capazes de fabricar mais de cerca de 11 g de efeito espúrio, e não exclui abaixo
disso. Onze gramas é da ordem do efeito procurado. Sem essa medida, o mesmo
placebo seria reportado como aprovação --- e teria sido.

\subsection{As ameaças que ficam declaradas}
\label{sub:ameacas-declaradas}

Três limitações não foram resolvidas e não devem ser apresentadas como se
tivessem sido.

A primeira é o \(d = 0\). O §2º do art.~28-B alcança a dispersão aérea para
controle vetorial, de modo que municípios sem agricultura pulverizada podem ser
tratados. A medida direta dessa contaminação existe no SINAN e soma 89
notificações em oito anos --- insuficiente para estimar o que quer que seja. A
única rota ao zero que mede o \emph{método} no pré-período é o cadastro estadual,
e ele depende de pedido de acesso à informação ainda pendente.

A segunda é a antecipação. A janela que define a dose começa em 2015, depois do
marco de notícia de fevereiro daquele ano, de modo que a dose medida pode já
responder à expectativa. A ameaça morde a variável de tratamento, não apenas o
desfecho, e a janela recuada de 2010--2014 entra como sensibilidade declarada
precisamente por isso.

A terceira é a hipótese que o alvo originalmente primário exigia. O teste de
inclinação em dose no pré-período não a rejeita, mas os três cortes placebo
produzem inclinação \emph{maior em módulo} que a do desenho real. O teste não
falha por o desenho ser limpo; ele não teria como rejeitar.

\subsection{O que vem depois}
\label{sub:depois}

Três frentes dependem de terceiros e têm relógio próprio: o pedido de acesso ao
cadastro estadual, que é a única via ao \(d = 0\) operacional; a verificação de
se o registro de nascimentos alcança 2023 e 2024, o que levaria o pós-banimento
de quatro para seis anos; e a varredura legislativa dos municípios restantes,
como conferência de completude.

\paragraph{E há uma hipótese que a literatura sugere e que este trabalho não
pode reivindicar.}
O efeito estimado por \citet{reynier2025glyphosate} concentra-se doze vezes mais
no decil inferior de peso esperado --- 75 g de perda ali contra 6 g no decil
superior. Confrontadas essas magnitudes com o efeito mínimo detectável deste
desenho, aparece uma assimetria que o diagnóstico agregado esconde: o desenho
não detecta o efeito médio, e detectaria com folga um efeito da ordem do que
aquele artigo encontra na cauda. Se o padrão distributivo se repetir aqui, o
problema deixa de ser de tamanho de amostra e passa a ser de escolha de
estimando.

\textbf{Isso não é resultado deste trabalho, e não pode ser lido como tal.} Não
foi declarado antes do dado, e promovê-lo agora seria exatamente a prática que a
pré-especificação existe para impedir. Registra-se como a hipótese a
pré-especificar no ciclo seguinte --- que é o único estatuto honesto que ela
pode ter aqui, e que vale mais do que um resultado a mais.

\subsection{Uma nota sobre o que este ensaio defende}
\label{sub:defende}

O trabalho não mostra que o banimento funcionou, e tampouco que não funcionou.
Mostra que o desenho disponível para avaliá-lo, montado com o estimador correto
para tratamento contínuo e submetido à camada de robustez que o estado da arte
recomenda, não distingue entre as duas coisas --- e diz exatamente quanto lhe
falta para distinguir.

Num campo em que a mesma combinação de registro administrativo amplo e poder
curto produz achados publicáveis com regularidade, estabelecer que um desenho
não fala, e estabelecê-lo contra critério fixado por escrito antes de olhar, é
contribuição de tipo diferente da que se buscava. Ela é menor em ambição e maior
em confiabilidade, e é a que o dado autoriza.
